%%%%%%%%%%%%%%%%%%%%%%%%%%%%%%%%%%%%%%%%%%%%%%%%%%%%%%%%%%%%%%%%%%%%%%%%%%%%%%%%
% Obxectivo: Lista de termos empregados no documento,                          %
%            xunto cos seus respectivos significados.                          %
%%%%%%%%%%%%%%%%%%%%%%%%%%%%%%%%%%%%%%%%%%%%%%%%%%%%%%%%%%%%%%%%%%%%%%%%%%%%%%%%

\newglossaryentry{udp}{
  name=UDP,
  description={Protocolo de comunicación de red ligero y rápido, usado para enviar paquetes sin establecer una conexión permanente entre emisor y receptor}
}
\newglossaryentry{sdk}{
  name=SDK,
  description={Kit de desarrollo de software (o por sus siglas en inglés, “Software Development Kit”) es un conjunto de 
  herramientas que permiten trabajar con una plataforma (por ejemplo NatNet con Motive y Optitrack)}
}
\newglossaryentry{lidar}{
  name=LiDAR,
  description={Sensor que utiliza la luz láser para medir las distancias y generar información sobre la geometría del entorno}
}
\newglossaryentry{deep learning}{
  name=Deep Learning,
  description={Rama del aprendizaje automático basada en redes neuronales profundas, capaz de aprender patrones complejos a partir de grandes conjuntos de datos}
}
\newglossaryentry{xml}{
  name=XML,
  description={Lenguaje de marcado utilizado para estructurar información en forma de etiquetas, empleado en este proyecto para definir los modelos de la simulación}
}
\newglossaryentry{sac}{
  name=SAC,
  description={Algoritmo de aprendizaje por refuerzo profundo (Sotf Actor-Critic), que combina aprendizaje off-policy con una política estocástica para favorecer la exploración}
}
\newglossaryentry{jac}{
  name=JAC,
  description={Herramientas empleadas para acelerar operaciones de cálculo durante el entrenamiento}
}
\newglossaryentry{jit}{
  name=JIT,
  description={Técnica de compilación \textit{Just-In-Time} que transforma partes del código en instrucciones optimizadas durante la ejecución para mejorar el rendimiento}
}
\newglossaryentry{mdp}{
  name=MDP,
  description={Proceso de decisión de Markov, modelo matemático utilizado para representar problemas de toma de decisiones mediante estados, acciones, recompensas y transiciones}
}
\newglossaryentry{off-policy}{
  name=Off-policy,
  description={Tipo de algoritmo de aprendizaje por refuerzo que puede aprender a partir de datos generados por una política distinta a la que se está optimizando}
}
\newglossaryentry{interpretado}{
  name=interpretado,
  description={Lenguaje de programación que se ejecuta directamente, sin necesidad de ser compilado previamente}
}
\newglossaryentry{api}{
  name=API,
  description={Conjunto de funciones, clases y métodos que permiten que distintos programas, librerías o sistemas se comuniquen entre sí. Facilita el uso de funcionalidades externas sin necesidad 
  de conocer su implementación interna}
}
\newglossaryentry{mjcf}{
  name=MJCF,
  description={MJCF (MuJoCo Modeling Format) es el formato de modelado XML usado en MuJoCo para describir modelos físicos}
}
\newglossaryentry{unicast}{
  name=Unicast,
  description={Modo de comunicación de red en el que los datos se envían desde un único emisor a un único receptor}
}
\newglossaryentry{imu}{
  name=IMU,
  description={Unidad de medición inercial formada por sensores como acelerómetros y giroscopios, utilidades para estimar la orientación, la aceleración y el movimiento de un dispositivo (en este 
  caso el robot Go2)}
}
\newglossaryentry{opengl}{
  name=OpenGL,
  description={Interfaz gráfica utilizada para renderizar gráficos en 2D y 3D, empleada habitualmente en simuladores y aplicaciones de visualización}
}
\newglossaryentry{markerset}{
  name=MarkerSet,
  description={Conjunto de marcadores reflectantes agrupados en Motive para identificar un mismo objeto dentro del entorno Optitrack}
}
\newglossaryentry{rigidbody}{
  name=RigidBody,
  description={Objeto definido en Motive a partir de varios marcadores, cuya posición y orientación pueden ser calculadas como un único cuerpo rígido}
}
\newglossaryentry{commit}{
  name=Commit,
  description={Registro de cambios realizado en un sistema de control de versiones, como GitHub, que guarda el estado de un proyecto en un momento determinado}
}
\newglossaryentry{ide}{
  name=IDE,
  description={Entorno de desarrollo integrado, permite escribir, ejecutar, depurar y gestionar código de manera cómoda}
}
\newglossaryentry{script}{
  name=script,
  description={Archivo de código que contiene una secuencia de instrucciones ejecutables}
}